% =============================================================================
% TCC -- Auditoria automática de laudos de mamografia
% Modelo abnTeX2 (ABNT NBR 14724). Compilar com: latexmk -pdf main.tex
% Marcadores: \pendente{...} aparece em vermelho no PDF e indica o que falta.
%             Para listar todos: procure "\pendente" nos arquivos .tex.
% =============================================================================
\documentclass[12pt, openright, oneside, a4paper, chapter=TITLE, english, brazil]{abntex2}

\usepackage[T1]{fontenc}
\usepackage[utf8]{inputenc}
\usepackage{lmodern}
\usepackage{indentfirst}
\usepackage{xcolor}
\usepackage{graphicx}
\usepackage{microtype}
\usepackage{booktabs}
\usepackage{multirow}
\usepackage{array}
\usepackage{icomma}          % vírgula decimal no modo matemático
\usepackage{amsmath}
\usepackage[brazilian, hyperpageref]{backref}
\usepackage[alf]{abntex2cite}

% ---------------------------------------------------------------- marcadores
\newcommand{\pendente}[1]{\textcolor{red}{[Pendente: #1]}}

% ---------------------------------------------------------------- backref
\renewcommand{\backrefpagesname}{Citado na(s) página(s):~}
\renewcommand{\backref}{}
\renewcommand*{\backrefalt}[4]{%
  \ifcase #1 Nenhuma citação no texto.%
  \or Citado na página #2.%
  \else Citado #1 vezes nas páginas #2.%
  \fi}

% ---------------------------------------------------------------- dados do trabalho
\titulo{Auditoria automática de laudos de mamografia: detecção de omissões e discordâncias com visão computacional e o léxico BI-RADS}
\autor{Erick Alves}
\local{Maceió}
\data{2026}
\orientador{\pendente{nome do orientador}}
\instituicao{%
  Universidade Federal de Alagoas
  \par
  Instituto de Computação
  \par
  \pendente{nome do curso}}
\tipotrabalho{Trabalho de Conclusão de Curso}
\preambulo{Trabalho de Conclusão de Curso apresentado ao curso de \pendente{nome do curso} do Instituto de Computação da Universidade Federal de Alagoas, como requisito parcial para a obtenção do grau de \pendente{grau}.}

\definecolor{blue}{RGB}{41,5,195}
\makeatletter
\hypersetup{
  pdftitle={\@title},
  pdfauthor={\@author},
  pdfsubject={\imprimirpreambulo},
  pdfcreator={LaTeX com abnTeX2},
  pdfkeywords={mamografia}{BI-RADS}{auditoria de laudos}{visão computacional}{detecção de objetos},
  colorlinks=true, linkcolor=blue, citecolor=blue, filecolor=magenta, urlcolor=blue,
  bookmarksdepth=4}
\makeatother

\setlength{\parindent}{1.3cm}
\setlength{\parskip}{0.2cm}
\makeindex

\begin{document}
\selectlanguage{brazil}
\frenchspacing

% ============================================================ PRÉ-TEXTUAIS
\imprimircapa
\imprimirfolhaderosto

\begin{folhadeaprovacao}
  \pendente{folha de aprovação: preencher com a banca depois da defesa}
\end{folhadeaprovacao}

\begin{agradecimentos}
  \pendente{agradecimentos}
\end{agradecimentos}

% Resumo e abstract: versão provisória. Reescrever no fim (S17), quando todos os
% resultados existirem. Limite ABNT: 150 a 500 palavras.

\setlength{\absparsep}{18pt}
\begin{resumo}
Laudos de mamografia apresentam inconsistências documentadas, como baixa adesão ao léxico BI-RADS e grafias divergentes de um mesmo termo, e não há ferramenta que confira automaticamente o que o laudo afirma contra o que a imagem mostra. Este trabalho propõe um sistema híbrido de auditoria que cruza a imagem mamográfica com o laudo já emitido e sinaliza dois tipos de problema: omissões de achados suspeitos e discordâncias de descritores BI-RADS. O sistema combina um detector de lesões treinado no VinDr-Mammo, um classificador de categoria BI-RADS e densidade por mama, um extrator de informação do laudo em português e um auditor baseado em regras rastreáveis. Na validação, o detector de massas alcançou AUC de 0,832 (IC 95\%: 0,768 a 0,889) em nível de mama e encontrou 61,2\% das massas com um falso positivo por imagem; uma ablação de pré-processamento não detectou diferença entre o recorte da mama, a ausência de recorte e o realce por CLAHE. \pendente{completar com os resultados do classificador, da extração dos laudos e da auditoria}

\textbf{Palavras-chave}: mamografia. BI-RADS. auditoria de laudos. visão computacional. detecção de objetos.
\end{resumo}

\begin{resumo}[Abstract]
\begin{otherlanguage*}{english}
\pendente{traduzir o resumo quando a versão final estiver pronta}

\vspace{\onelineskip}
\noindent
\textbf{Keywords}: mammography. BI-RADS. report auditing. computer vision. object detection.
\end{otherlanguage*}
\end{resumo}


\pdfbookmark[0]{\listfigurename}{lof}
\listoffigures*
\cleardoublepage
\pdfbookmark[0]{\listtablename}{lot}
\listoftables*
\cleardoublepage

\begin{siglas}
  \item[ACR] \textit{American College of Radiology}
  \item[AP] \textit{Average Precision} (precisão média)
  \item[AUC] Área sob a curva ROC
  \item[BI-RADS] \textit{Breast Imaging Reporting and Data System}
  \item[CC] Vista craniocaudal
  \item[CLAHE] \textit{Contrast Limited Adaptive Histogram Equalization}
  \item[DICOM] \textit{Digital Imaging and Communications in Medicine}
  \item[FROC] \textit{Free-response Receiver Operating Characteristic}
  \item[IC] Intervalo de confiança
  \item[LLM] \textit{Large Language Model} (modelo de linguagem de grande porte)
  \item[MLO] Vista mediolateral oblíqua
  \item[ROI] Região de interesse
  \item[VLM] \textit{Vision-Language Model} (modelo de visão e linguagem)
\end{siglas}

\pdfbookmark[0]{\contentsname}{toc}
\tableofcontents*
\cleardoublepage

% ============================================================ TEXTUAIS
\textual
% =============================================================================
\chapter{Introdução}\label{cap:introducao}
% =============================================================================

\section{Contextualização}

A mamografia é o exame de imagem usado no rastreamento do câncer de mama, e o seu resultado chega ao médico solicitante e à paciente na forma de um laudo em texto. Para tornar esses laudos comparáveis e menos ambíguos, o \textit{American College of Radiology} mantém o BI-RADS, um léxico padronizado que define como descrever os achados (massas, calcificações, distorções e assimetrias), como classificar a composição da mama e como atribuir uma categoria final de avaliação que orienta a conduta \cite{dorsi2013birads}. \pendente{incluir um dado oficial de incidência do câncer de mama no Brasil (INCA), com referência}

Na prática, a adesão ao léxico é incompleta. Ao analisar 22.247 laudos eletrônicos de mamografia de um hospital universitário brasileiro, \citeonline{serapiao2010} encontraram que apenas 11\% a 61\% dos laudos continham os termos BI-RADS formatados corretamente, que 21\% dos termos descritivos tinham erro de grafia e que a palavra ``microcalcificação'' aparecia escrita de 36 formas diferentes. Esses números mostram que o laudo, que é o registro oficial do exame, pode divergir do padrão e, por consequência, da própria imagem que descreve.

\section{Problema}

Hoje, a conferência entre o que o laudo afirma e o que a imagem mostra, quando acontece, é manual. Não há, ao conhecimento do autor, ferramenta que faça essa verificação de forma automática em mamografia. O problema tratado neste trabalho é, portanto, o seguinte: \textbf{dado um exame de mamografia e o laudo já emitido para ele, identificar automaticamente (i) achados suspeitos presentes na imagem que o laudo não menciona e (ii) descrições do laudo que não conferem com a imagem.}

O primeiro caso é chamado aqui de \textit{omissão}; o segundo, de \textit{discordância}. Em ambos, a saída esperada não é um diagnóstico, e sim um alerta acompanhado da evidência que o motivou, para que o radiologista decida se revisa o laudo.

\section{Objetivos}

\subsection{Objetivo geral}

Desenvolver e avaliar um sistema de auditoria automática de laudos de mamografia que cruze a imagem com o laudo e sinalize omissões de achados suspeitos e discordâncias de descritores BI-RADS, com alertas rastreáveis.

\subsection{Objetivos específicos}

\begin{enumerate}
  \item Construir um pipeline de pré-processamento de mamografias em formato DICOM, com recorte automático da região da mama, e avaliar o efeito das escolhas de pré-processamento por meio de uma ablação.
  \item Treinar e avaliar um detector de lesões suspeitas, medindo o desempenho em nível de lesão e em nível de mama, com intervalos de confiança.
  \item Treinar e avaliar um classificador de categoria BI-RADS e de densidade mamária por mama.
  \item Desenvolver um extrator que converta laudos em português em uma representação estruturada compatível com o BI-RADS.
  \item Definir um auditor baseado em regras explícitas que compare as saídas da imagem e do laudo e gere alertas com severidade e evidência.
  \item Avaliar o sistema completo em um conjunto de laudos reais com erros injetados de forma controlada e comparar com alternativas de referência.
\end{enumerate}

\section{Contribuições}

As contribuições pretendidas deste trabalho são:

\begin{itemize}
  \item a aplicação da auditoria automática entre imagem e laudo à mamografia, tarefa até então estudada em radiografia de tórax \cite{mahmood2025factcheck, zou2025corbenchx};
  \item o uso do léxico BI-RADS como esquema de verificação, e não apenas como rótulo de treinamento;
  \item o tratamento de laudos em português, usando os laudos reais distribuídos com a base INbreast \cite{moreira2012inbreast};
  \item uma ablação controlada de pré-processamento para detecção de massas, que separa o efeito do recorte da mama e o do realce de contraste;
  \item uma avaliação com separação estrita entre validação e teste e com intervalos de confiança por reamostragem de exames.
\end{itemize}

\section{Escopo e limitações}

O escopo foi delimitado por características dos dados e pelo prazo do trabalho. A base VinDr-Mammo só marca com caixa os achados de categoria BI-RADS 3 ou superior \cite{nguyen2023vindr}; por isso, ``omissão'' significa aqui omissão de achado \textit{suspeito}. O caminho de detecção de calcificações foi retirado do escopo, e o detector concentra-se em massas e em outros achados suspeitos de maior tamanho. O trabalho não contou com revisão por radiologista: as anotações de laudo foram produzidas por anotadores treinados em um protocolo escrito, e a referência de achados em imagem vem das anotações públicas das bases. Os resultados descrevem, portanto, a consistência do sistema frente a esses gabaritos, e não a sua validade diagnóstica em uso clínico.

\section{Organização do trabalho}

O \autoref{cap:fundamentacao} apresenta os conceitos de mamografia, BI-RADS, detecção de objetos e avaliação de detectores, além dos trabalhos relacionados. O \autoref{cap:metodologia} descreve as bases de dados, o pré-processamento, os modelos, o auditor e o protocolo de avaliação. O \autoref{cap:resultados} reporta os resultados, discutidos no \autoref{cap:discussao}. O \autoref{cap:conclusao} conclui o trabalho e aponta trabalhos futuros.

% =============================================================================
\chapter{Fundamentação teórica e trabalhos relacionados}\label{cap:fundamentacao}
% =============================================================================

\section{Mamografia e o léxico BI-RADS}

O exame de mamografia de rastreamento registra, em geral, duas vistas de cada mama: a craniocaudal (CC) e a mediolateral oblíqua (MLO). Uma mesma lesão costuma aparecer nas duas vistas, em posições diferentes, o que permite ao radiologista confirmar um achado e estimar a sua localização.

O BI-RADS organiza o laudo em três partes principais \cite{dorsi2013birads}. A primeira é a composição da mama, em quatro categorias de densidade (A a D), que vão da mama predominantemente adiposa à extremamente densa. A segunda é a descrição dos achados com um vocabulário controlado: massas, descritas por forma, margem e densidade; calcificações, por morfologia e distribuição; distorções arquiteturais; e assimetrias. A terceira é a categoria final de avaliação, de 0 a 6, em que as categorias 4 e 5 indicam suspeita de malignidade e recomendação de biópsia. Essa estrutura é o que torna possível comparar automaticamente um laudo com a imagem: cada termo do laudo corresponde a um conceito definido, que pode ser buscado na imagem.

\section{Detecção de objetos com redes neurais}

Um detector de objetos recebe uma imagem e devolve um conjunto de caixas delimitadoras, cada uma com uma classe e uma confiança. Os detectores de estágio único, como a família YOLO, fazem a localização e a classificação em uma única passagem pela rede, o que os torna rápidos e adequados a recursos computacionais limitados \cite{jocher2023yolov8}. Neste trabalho, o detector é usado para localizar achados suspeitos na mamografia, e a confiança de cada caixa é a base para decidir se um achado deve gerar alerta.

\pendente{descrever brevemente a arquitetura do YOLOv8 (backbone, neck, cabeça sem âncoras) e incluir uma figura}

\section{Avaliação de detectores em imagem médica}

\subsection{Interseção sobre união e precisão média}

Uma detecção é considerada correta quando a sua sobreposição com uma caixa de referência, medida pela razão entre a interseção e a união das duas áreas (IoU), passa de um limiar, tipicamente 0,5. A precisão média (AP) resume a curva de precisão e revocação de uma classe, e o mAP é a média do AP entre as classes. O AP é a métrica padrão em visão computacional, mas é pouco sensível ao número absoluto de falsos positivos por imagem, que é o que determina a carga de trabalho de quem revisa os alertas \cite{maierhein2024metrics}.

\subsection{Análise FROC}

A análise FROC (\textit{Free-response Receiver Operating Characteristic}) relaciona a sensibilidade em nível de lesão com o número médio de falsos positivos por imagem, à medida que o limiar de confiança varia \cite{bunch1978froc}. Diferente do AP, o eixo horizontal da FROC tem interpretação direta: ``encontrar 60\% das lesões com um falso positivo por imagem'' é uma afirmação que um radiologista consegue avaliar. Por isso, a FROC é a métrica tradicional da detecção assistida por computador em mamografia. Para resumir a curva em um número, usa-se a área parcial sob a FROC em uma faixa clinicamente relevante de falsos positivos por imagem, normalizada para o intervalo de 0 a 1.

\subsection{Intervalos de confiança por reamostragem}

Conjuntos de validação em imagem médica costumam ser pequenos, e uma diferença de poucos pontos entre dois modelos pode ser apenas ruído. O \textit{bootstrap} estima a variabilidade de uma métrica reamostrando com reposição as unidades do conjunto de avaliação \cite{efron1993bootstrap}. Em mamografia, a unidade de reamostragem deve ser o exame, e não a imagem, porque as vistas de um mesmo exame são correlacionadas; reamostrar imagens isoladas subestima a variância.

\section{Trabalhos relacionados}

\subsection{Auditoria entre imagem e laudo}

A verificação automática de laudos contra a imagem foi estudada principalmente em radiografia de tórax. \citeonline{mahmood2025factcheck} propõem um modelo de checagem de fatos ancorado em frases, que detecta erros nos achados e nas suas localizações em laudos gerados automaticamente; a detecção de achados omitidos não é tratada no trabalho. \citeonline{zou2025corbenchx} constroem o CorBenchX, um conjunto de laudos de tórax com erros e um \textit{benchmark} de correção por modelos de visão e linguagem, no qual o melhor modelo comercial avaliado detectou corretamente 50,6\% dos erros. \pendente{incluir RADAR (2026) e MammoVQA, com referências conferidas}

Esses resultados indicam duas lacunas. A primeira é de modalidade: não foi encontrado trabalho que audite laudos de mamografia contra a imagem, nem que use o léxico BI-RADS como esquema de verificação. A segunda é de abordagem: o desempenho modesto dos modelos genéricos de visão e linguagem justifica investigar sistemas com módulos especializados.

\subsection{Detecção e classificação em mamografia}

O VinDr-Mammo é a maior base pública de mamografia digital com anotações de achados por caixa e categorias BI-RADS por mama e por achado, com divisão oficial entre treino e teste \cite{nguyen2023vindr}. O CBIS-DDSM reúne mamografias de filme digitalizado com máscaras de lesões \cite{lee2017cbis}, e o INbreast contém mamografias digitais com anotações detalhadas \cite{moreira2012inbreast}. Entre os modelos recentes, o Mammo-CLIP, um modelo de visão e linguagem pré-treinado em mamografias, relata mAP de cerca de 0,58 na detecção de massas no VinDr-Mammo, e destaca como principal vantagem a eficiência no uso de rótulos \cite{ghosh2024mammoclip}.

\subsection{Pré-processamento}

Duas etapas de pré-processamento são frequentes em mamografia. O recorte da região da mama remove o fundo e aumenta a resolução efetiva da mama na entrada da rede; em classificação no VinDr-Mammo, o recorte elevou a área sob a curva ROC de 79,6 para 83,8 \cite{ibragimov2024mamt4}. O realce de contraste por CLAHE \cite{zuiderveld1994clahe} também é comum, mas a evidência de ganho em aprendizado profundo para mamografia é fraca, e muitos estudos aplicam o recorte e o realce juntos, sem separar o efeito de cada um. \pendente{citar os estudos de ablação de CLAHE (Cancers 16(2):322; Algorithms 18(12):796)} Essa lacuna motiva a ablação apresentada neste trabalho.

\subsection{Laudos em português}

Do lado do texto, \citeonline{serapiao2010} documentaram a baixa adesão ao léxico BI-RADS em laudos brasileiros, usando mineração de texto. \pendente{incluir trabalhos de extração de informação de laudos com LLM e com regras, e recursos em português (BioBERTpt, BRAX)}

% =============================================================================
\chapter{Materiais e métodos}\label{cap:metodologia}
% =============================================================================

\section{Visão geral do sistema}

O sistema proposto tem quatro componentes (\autoref{fig:pipeline}). Do lado da imagem, um \textbf{detector de lesões} localiza achados suspeitos em cada vista, e um \textbf{classificador por mama} estima a categoria BI-RADS e a densidade. Do lado do texto, um \textbf{extrator} converte o laudo em uma representação estruturada. Por fim, um \textbf{auditor} baseado em regras compara as duas representações e gera alertas com severidade e evidência. A escolha por módulos especializados, e não por um único modelo de visão e linguagem, segue o desempenho modesto desses modelos na tarefa de verificar laudos (\autoref{cap:fundamentacao}) e privilegia a rastreabilidade: cada alerta pode ser explicado por uma regra e por um limiar declarado.

\begin{figure}[htb]
  \centering
  \fbox{\parbox{0.85\textwidth}{\centering\vspace{1.5cm}\pendente{figura do pipeline: imagem $\rightarrow$ detector e classificador; laudo $\rightarrow$ extrator; ambos $\rightarrow$ auditor $\rightarrow$ alerta}\vspace{1.5cm}}}
  \caption{Visão geral do sistema de auditoria.}
  \label{fig:pipeline}
  \fonte{Elaborada pelo autor.}
\end{figure}

\section{Bases de dados}

Foram usadas três bases públicas, com papéis distintos (\autoref{tab:bases}).

\begin{table}[htb]
  \centering
  \caption{Bases de dados usadas e o papel de cada uma no trabalho.}
  \label{tab:bases}
  \small
  \begin{tabular}{@{}p{2.4cm}p{2.6cm}p{4.6cm}p{4.2cm}@{}}
    \toprule
    Base & Tipo de imagem & Conteúdo usado & Papel no trabalho \\
    \midrule
    VinDr-Mammo \cite{nguyen2023vindr} & Digital (FFDM) & 5.000 exames, 20.000 imagens; caixas de achados com BI-RADS 3 a 5; BI-RADS e densidade por mama & Treino, validação e teste do módulo de imagem \\
    CBIS-DDSM \cite{lee2017cbis} & Filme digitalizado & Massas com máscara de segmentação, convertidas em caixas & Teste externo \\
    INbreast \cite{moreira2012inbreast} & Digital (FFDM) & 410 imagens e 117 laudos em português & Corpus de laudos da auditoria \\
    \bottomrule
  \end{tabular}
  \fonte{Elaborada pelo autor.}
\end{table}

No VinDr-Mammo, o arquivo de anotação de achados tem 20.486 linhas, das quais 2.254 têm caixa delimitadora. As caixas existem apenas para achados de categoria BI-RADS 3, 4 e 5; achados benignos não são marcados. A categoria BI-RADS por mama varia de 1 a 5 (não há categoria 0), e a densidade A ocorre em apenas 50 mamas, motivo pelo qual as densidades A e B foram fundidas em uma única classe. A distribuição por mama é mostrada na \autoref{tab:dist_mamas}.

No CBIS-DDSM, as caixas de massa foram derivadas das máscaras de segmentação, porque as colunas de caminho dos arquivos de metadados trocam de forma inconsistente a imagem recortada e a máscara; cada par foi resolvido pelo conteúdo do arquivo. A conversão do subconjunto de massas resultou em 1.592 imagens e 1.696 caixas. Pacientes presentes tanto no treino quanto no teste oficiais (31) foram movidos inteiros para o teste.

Os laudos do INbreast são arquivos de texto em português, codificados em Latin-1, que contêm a categoria BI-RADS explícita, descritores, quadrantes e negações. A literatura que descreve a base não documenta esses laudos \cite{moreira2012inbreast}; por isso, eles são descritos aqui a partir da inspeção direta do pacote de dados. Por cuidado ético, nenhum laudo é reproduzido integralmente neste texto.

\section{Divisão dos dados}

A divisão segue duas regras. A primeira é usar o teste oficial do VinDr-Mammo (1.000 exames) como teste, sem alterá-lo, para que os resultados sejam comparáveis à literatura. A segunda é que nenhuma decisão de modelagem pode olhar o teste. Para isso, foi criado um conjunto de validação com cerca de 10\% dos exames do treino oficial. A escolha é feita por exame, de modo que as quatro vistas de um exame fiquem sempre no mesmo conjunto, e é determinística: um exame vai para a validação quando o valor derivado do resumo SHA-256 do seu identificador, somado a uma constante fixa, é menor que 0,10. A mesma validação é usada pelo detector e pelo classificador, o que permite calibrar o auditor em um único conjunto sem vazamento entre os módulos.

\begin{table}[htb]
  \centering
  \caption{Número de mamas por categoria BI-RADS e por densidade em cada conjunto (VinDr-Mammo).}
  \label{tab:dist_mamas}
  \small
  \begin{tabular}{@{}lrrrrrrrrr@{}}
    \toprule
    & \multicolumn{5}{c}{BI-RADS} & \multicolumn{3}{c}{Densidade} & \\
    \cmidrule(lr){2-6}\cmidrule(lr){7-9}
    Conjunto & 1 & 2 & 3 & 4 & 5 & A+B & C & D & Total \\
    \midrule
    Treino    & 4.821 & 1.677 & 338 & 276 & 80 & 724 & 5.494 & 974 & 7.192 \\
    Validação &   541 &   194 &  34 &  29 & 10 &  80 &   622 & 106 &   808 \\
    Teste     & 1.341 &   467 &  93 &  76 & 23 & 200 & 1.530 & 270 & 2.000 \\
    \bottomrule
  \end{tabular}
  \fonte{Elaborada pelo autor a partir de \citeonline{nguyen2023vindr}.}
\end{table}

Para o detector, cada imagem recebeu um papel. Imagens com ao menos uma caixa de classe do detector são \textit{positivas}; imagens sem nenhuma caixa são \textit{normais}; e 189 imagens cujas caixas eram todas de calcificação ficaram de fora, por não serem nem positivas nem normais limpas para um detector de massas. Validação e teste recebem todas as imagens normais, para medir falsos positivos com a prevalência real. O treino recebe uma imagem normal para cada imagem positiva, escolhidas de forma determinística. A \autoref{tab:split_detector} resume o resultado.

\begin{table}[htb]
  \centering
  \caption{Divisão das imagens do VinDr-Mammo para o detector.}
  \label{tab:split_detector}
  \begin{tabular}{@{}lrrrr@{}}
    \toprule
    Conjunto & Imagens positivas & Imagens normais & Caixas & Exames \\
    \midrule
    Treino    & 1.148 & 1.148 & 1.358 & 1.514 \\
    Validação &   121 & 1.484 &   160 &   403 \\
    Teste     &   310 & 3.643 &   367 & 1.000 \\
    \bottomrule
  \end{tabular}
  \fonte{Elaborada pelo autor.}
\end{table}

\section{Pré-processamento}

O pré-processamento parte do arquivo DICOM e segue a ordem prevista no padrão: aplicação da transformação de modalidade (\textit{rescale}), recorte da mama, ajuste de intensidade, correção de polaridade, espelhamento canônico e redimensionamento. Imagens marcadas como ``FOR PROCESSING'' são descartadas.

\textbf{Recorte da mama.} A região da mama é localizada por limiarização de Otsu \cite{otsu1979} sobre uma cópia da imagem em que os pixels mais claros foram suprimidos, para remover etiquetas, marcadores metálicos e bordas do detector. O nível de supressão é o percentil 95 calculado apenas sobre os pixels que não são fundo. Esse detalhe é necessário: quando o percentil era calculado sobre o quadro inteiro, em mamas pequenas (que ocupam de 6\% a 10\% do quadro) ele caía dentro da própria mama, o tecido denso era apagado e o recorte era abandonado. Após a limiarização, a imagem é suavizada e dilatada, e a caixa envolvente do maior contorno define o recorte. Contornos com menos de 4\% da área do quadro são rejeitados.

\textbf{Intensidade e polaridade.} A intensidade é ajustada pela janela (centro e largura) e pela função de transformação registradas no cabeçalho DICOM; na ausência delas, usa-se uma normalização pelos percentis 1 e 99 dentro da máscara da mama. Imagens com interpretação fotométrica MONOCHROME1 são invertidas depois do janelamento, para que o tecido apareça sempre claro sobre fundo escuro.

\textbf{Orientação e tamanho.} As imagens da mama direita são espelhadas horizontalmente, de modo que a parede torácica fique sempre à esquerda. Em seguida, a imagem é redimensionada preservando a proporção e completada com zeros até 1.024 $\times$ 640 pixels (altura $\times$ largura). As caixas de referência passam pelas mesmas transformações (deslocamento do recorte, espelhamento, escala e preenchimento), o que é verificado por testes automáticos.

\textbf{Braços da ablação.} Para medir o efeito das escolhas de pré-processamento, foram definidos três braços (\autoref{tab:bracos}), todos com o mesmo ajuste de intensidade pela janela DICOM. O braço B2 isola o efeito do realce de contraste, e o B4, o do recorte.

\begin{table}[htb]
  \centering
  \caption{Braços da ablação de pré-processamento.}
  \label{tab:bracos}
  \begin{tabular}{@{}llp{6.5cm}@{}}
    \toprule
    Braço & Recorte da mama & Realce de contraste \\
    \midrule
    B0 (referência) & sim & nenhum \\
    B2 & sim & CLAHE \cite{zuiderveld1994clahe}, com limite de 2,0 e blocos de 8 $\times$ 8 \\
    B4 & não & nenhum \\
    \bottomrule
  \end{tabular}
  \fonte{Elaborada pelo autor.}
\end{table}

\section{Detector de lesões}

O detector é um YOLOv8s \cite{jocher2023yolov8} inicializado com pesos pré-treinados no COCO e treinado para seis classes: massa, distorção arquitetural, assimetria focal, assimetria, assimetria global e ``outro achado suspeito'', que agrupa espessamento e retração de pele, retração de mamilo e linfonodo suspeito, todos com poucos exemplos. A entrada tem lado maior de 1.024 pixels, com lotes retangulares que respeitam a proporção da imagem.

O treino usa o otimizador AdamW com taxa de aprendizado inicial de 0,001 e decaimento em cosseno, aquecimento de três épocas, lote de duas imagens com acumulação de gradiente até um lote nominal de 16, precisão mista e no máximo 80 épocas, com parada antecipada após 15 épocas sem melhora na validação. O aumento de dados é geométrico e leve (espelhamento horizontal, rotação de até 15°, translação de 5\% e escala de 15\%) e inclui variação de brilho; não há espelhamento vertical, \textit{mixup} nem cópia e colagem de objetos. A semente aleatória é fixa e o treino é determinístico.

\pendente{decidir com o orientador se as classes raras serão fundidas em uma classe ou se o detector ficará só com massa, e atualizar esta seção}

\section{Classificador por mama}

O classificador estima a categoria BI-RADS (cinco classes) e a densidade (três classes) de cada mama. A arquitetura é uma EfficientNet-B2 \cite{tan2019efficientnet} pré-treinada no ImageNet, com entrada de um canal e duas cabeças lineares sobre as mesmas características. Como o BI-RADS é ordinal, a perda da categoria é a distância do transportador (\textit{Earth Mover's Distance}) ao quadrado entre as distribuições acumuladas prevista e real \cite{hou2016emd}, que penaliza mais os erros distantes; a densidade usa entropia cruzada com peso 0,5.

O modelo é treinado por imagem, com o rótulo da mama, e a previsão por mama é a média das probabilidades das suas vistas. Para compensar o desequilíbrio das classes (\autoref{tab:dist_mamas}), as imagens são sorteadas com peso proporcional à raiz do inverso da frequência da sua categoria BI-RADS. O treino usa AdamW (taxa de 0,0002, decaimento de pesos de 0,0001), decaimento em cosseno, lote de quatro imagens com acumulação até 16, precisão mista e parada antecipada pelo kappa ponderado na validação. O aumento de dados inclui rotação de até 10°, translação, escala e variação de brilho e contraste, sem espelhamento, para preservar a orientação canônica.

\pendente{incluir a comparação com o Mammo-CLIP \cite{ghosh2024mammoclip} se for aprovada}

\section{Extração de informação do laudo}

O laudo é convertido em uma representação estruturada comum aos dois lados do sistema, em que cada achado tem categoria, lateralidade, região, descritores BI-RADS e um trecho do texto como evidência. Cada campo distingue três estados: \textit{afirmado}, \textit{negado} e \textit{ausente}. A distinção entre ``o laudo diz que não há'' e ``o laudo não diz'' é essencial para a auditoria, porque só o primeiro caso pode contradizer a imagem.

O extrator de referência (E0) é baseado em regras: o laudo é dividido em seções, as seções que não se referem à mamografia (por exemplo, ultrassonografia e axila) são descartadas, o texto é normalizado (minúsculas e sem acentos) e os termos são buscados em um léxico BI-RADS bilíngue, com tratamento de negação pelo contexto da frase. A categoria BI-RADS escrita no laudo, incluindo as subcategorias 4A, 4B e 4C, é extraída por expressão regular.

\pendente{descrever o extrator com LLM local (E1): modelo, quantização, saída em JSON restrita por gramática e prompts}

\section{Auditor}

O auditor é determinístico. Ele compara a representação da imagem com a do laudo e aplica as regras da \autoref{tab:regras}. Duas regiões são consideradas compatíveis quando são iguais, adjacentes ou quando uma delas não foi especificada. Cada alerta recebe uma severidade de 1 (sem erro) a 5 (emergente) \pendente{citar a origem da escala de severidade} e carrega a evidência que o motivou: a caixa e a confiança da imagem e o trecho do laudo.

\begin{table}[htb]
  \centering
  \caption{Regras do auditor.}
  \label{tab:regras}
  \small
  \begin{tabular}{@{}lp{4.3cm}p{8cm}@{}}
    \toprule
    Código & Tipo & Condição \\
    \midrule
    A1 & Omissão de achado suspeito & A imagem contém um achado com confiança acima do limiar e o laudo não o afirma. \\
    A2 & Achado sem sustentação & O laudo afirma um achado que a imagem não ancora. \\
    A3 & Discordância de lateralidade & O achado está na mama oposta à indicada no laudo. \\
    A4 & Discordância de localização & A região do laudo não é compatível com a da imagem. \\
    A5 & Discordância de descritor & Forma, margem ou outro descritor BI-RADS diverge. \\
    A6 & Inconsistência de categoria & A categoria BI-RADS do laudo diverge da estimada. \\
    A7 & Discordância de densidade & A densidade do laudo diverge da estimada. \\
    A8 & Inconsistência interna & O próprio laudo é contraditório, sem depender da imagem. \\
    \bottomrule
  \end{tabular}
  \fonte{Elaborada pelo autor.}
\end{table}

O limiar de confiança que transforma uma detecção em achado (usado na regra A1) é calibrado na validação, no ponto de operação escolhido da curva FROC, e congelado antes de qualquer avaliação no teste. \pendente{registrar o ponto de operação escolhido e o valor do limiar}

\section{Protocolo de avaliação}

Todas as comparações entre variantes e todas as escolhas de limiar são feitas na validação. O teste oficial é usado uma única vez, com os modelos congelados, e o CBIS-DDSM é usado como teste externo.

\textbf{Detector.} São reportados: (i) o AP com IoU de 0,5 por classe, na validação completa e apenas nas imagens com achado, para separar o efeito das imagens normais; (ii) a FROC em nível de lesão, com acerto quando a IoU com a caixa de referência é de pelo menos 0,5, e a sensibilidade em 0,125, 0,25, 0,5, 1, 2 e 4 falsos positivos por imagem; (iii) a área parcial sob a FROC entre 0,125 e 4 falsos positivos por imagem, normalizada; e (iv) métricas por mama, em que a pontuação da mama é a maior confiança de massa entre as suas vistas e a mama é positiva quando alguma vista tem caixa de massa. Por mama, reportam-se a AUC e a sensibilidade quando 10\% e 20\% das mamas normais geram alerta. As predições são geradas com confiança mínima de 0,001 e até 20 caixas por imagem, para que a curva FROC alcance 4 falsos positivos por imagem.

\textbf{Intervalos de confiança.} Todos os intervalos de confiança de 95\% são obtidos por \textit{bootstrap} percentil com 1.000 reamostragens de exames \cite{efron1993bootstrap}, com semente fixa. Quando os intervalos de duas variantes se sobrepõem, o texto afirma que não foi possível detectar diferença, e não que as variantes são equivalentes.

\textbf{Classificador.} A métrica principal é o kappa ponderado quadrático da categoria BI-RADS por mama, adequado a classes ordinais \cite{cohen1960kappa}. Também são reportados a macro-F1, a acurácia, a AUC para BI-RADS 4 ou superior (a fronteira que indica biópsia) e, para a densidade, a acurácia e o kappa ponderado.

\textbf{Extração e auditoria.} \pendente{descrever: anotação de cerca de 110 laudos do INbreast com manual de anotação, 30 laudos em duplicata e kappa de Cohen \cite{cohen1960kappa}; métricas por campo do extrator; benchmark de auditoria com erros injetados em laudos reais, condicionados ao contexto; validação de plausibilidade por dois avaliadores; comparação com as referências Z0 (laudo sempre correto), Z1 (léxica) e Z2 (modelo de visão e linguagem sem treino)}

\section{Ambiente computacional}

Os experimentos foram executados em um computador pessoal com placa de vídeo NVIDIA GeForce RTX 3070 (8 GB), em Python 3.11, com PyTorch 2.13 (CUDA 12.6), Ultralytics 8.4 e timm 1.0. O código está em um repositório versionado e é coberto por testes automáticos, que verificam, entre outros pontos, a ausência de vazamento entre conjuntos, o mapeamento das caixas no pré-processamento e o cálculo das métricas. Os dados e os pesos dos modelos não são versionados, em respeito às licenças das bases.

\section{Aspectos éticos e anotação}

As bases usadas são públicas e anonimizadas, e foram obtidas conforme os termos de uso de cada uma. O trabalho não contou com revisão por radiologista. A anotação dos laudos é uma tarefa de extração do que já está escrito, e não de julgamento clínico, e foi feita por anotadores treinados em um manual de anotação escrito; a concordância entre eles mede a consistência do esquema de anotação, e não acurácia clínica. A análise de erro limita-se a causas técnicas observáveis, com o gabarito das bases como referência.

% =============================================================================
\chapter{Resultados}\label{cap:resultados}
% =============================================================================

Todos os resultados deste capítulo, exceto os da \autoref{sec:res_teste}, foram obtidos no conjunto de validação. Os intervalos entre colchetes são intervalos de confiança de 95\% por \textit{bootstrap} de exames.

\section{Verificação do pré-processamento}

Uma inspeção visual de 300 imagens processadas mostrou 20 casos em que o recorte da mama falhou e a imagem foi mantida sem recorte. Todos eram vistas craniocaudais de mamas pequenas, que ocupavam de 6\% a 10\% do quadro. A causa foi o nível de supressão de pixels claros, calculado sobre o quadro inteiro: como o fundo dominava a imagem, o percentil 95 caía dentro da mama e apagava o tecido denso. Após calcular o percentil apenas sobre os pixels que não são fundo, como descrito no \autoref{cap:metodologia}, a área recortada nesses casos passou a coincidir com a da mama. A correção foi aplicada antes de todos os treinos reportados a seguir.

\pendente{incluir figura com um exemplo antes e depois da correção do recorte}

\section{Detector de lesões}

\subsection{Escolha da configuração de treino}

A \autoref{tab:runs_detector} compara três treinos do detector no braço B0. A métrica é o AP com IoU de 0,5 para massa, medido na validação completa (com todas as imagens normais) e apenas nas imagens com achado. A diferença entre as duas últimas colunas mostra o peso das imagens normais: o mesmo modelo perde mais da metade do AP quando as 1.484 imagens normais da validação entram na conta, porque cada falso positivo nelas reduz a precisão.

\begin{table}[htb]
  \centering
  \caption{AP de massa (IoU de 0,5) na validação para três configurações de treino do detector.}
  \label{tab:runs_detector}
  \small
  \begin{tabular}{@{}lcccc@{}}
    \toprule
    Resolução & Normais no treino & Normais por positiva & Validação completa & Só com achado \\
    \midrule
    1.024 & 287 & 0,25 & 0,171 & 0,428 \\
    \textbf{1.024} & \textbf{1.148} & \textbf{1} & \textbf{0,186} & \textbf{0,448} \\
    1.536 & 1.148 & 1 & 0,197 & 0,384 \\
    \bottomrule
  \end{tabular}
  \fonte{Elaborada pelo autor.}
\end{table}

Aumentar o número de imagens normais no treino melhorou as duas colunas. A resolução de 1.536 pixels aumentou o AP na validação completa, mas reduziu o AP nas imagens com achado, ou seja, gerou menos falsos positivos em imagens normais à custa de encontrar menos lesões. Como as diferenças são pequenas diante do tamanho da validação e o treino em 1.536 pixels é mais caro, a configuração escolhida e congelada foi a de 1.024 pixels com uma normal para cada positiva (melhor época: 36).

As classes raras do detector (distorção arquitetural, os três tipos de assimetria e ``outro achado suspeito'') têm de 2 a 21 instâncias na validação, e o AP de todas ficou próximo de zero. Com tão poucos exemplos, não é possível estimar o desempenho nessas classes, e os resultados seguintes tratam apenas de massa.

\subsection{Desempenho em nível de lesão e de mama}

A validação tem 103 caixas de massa e, em nível de mama, 50 mamas com massa e 753 sem. A \autoref{tab:froc_b0} mostra a sensibilidade em nível de lesão do modelo congelado nos pontos de operação da curva FROC.

\begin{table}[htb]
  \centering
  \caption{Sensibilidade em nível de lesão (massa) por número de falsos positivos por imagem (FPPI), braço B0, validação.}
  \label{tab:froc_b0}
  \begin{tabular}{@{}lcccccc@{}}
    \toprule
    FPPI & 0,125 & 0,25 & 0,5 & 1 & 2 & 4 \\
    \midrule
    Sensibilidade & 0,408 & 0,466 & 0,563 & 0,612 & 0,689 & 0,786 \\
    \bottomrule
  \end{tabular}
  \fonte{Elaborada pelo autor.}
\end{table}

A área parcial normalizada sob a FROC, entre 0,125 e 4 FPPI, foi de 0,680 [0,577 a 0,782]. Em nível de mama, a AUC foi de 0,832 [0,768 a 0,889]. Quando o limiar é ajustado para que 10\% das mamas normais gerem alerta, o detector encontra 56,0\% [40,7\% a 70,0\%] das mamas com massa; com 20\% de alertas em mamas normais, encontra 72,0\% [57,1\% a 84,2\%].

\pendente{incluir a figura da curva FROC com a faixa de confiança}

\subsection{Ablação de pré-processamento}

Os braços B2 (com CLAHE) e B4 (sem recorte) foram treinados com a mesma configuração e avaliados da mesma forma que o B0. A \autoref{tab:ablacao} resume o resultado.

\begin{table}[htb]
  \centering
  \caption{Ablação de pré-processamento para detecção de massas na validação.}
  \label{tab:ablacao}
  \footnotesize
  \begin{tabular}{@{}lccc@{}}
    \toprule
    Métrica & B0 (referência) & B2 (+CLAHE) & B4 (sem recorte) \\
    \midrule
    Sensibilidade a 0,5 FPPI & 0,563 & 0,485 & 0,563 \\
    Sensibilidade a 1 FPPI   & 0,612 & 0,563 & 0,650 \\
    pAUC FROC (0,125 a 4)    & 0,680 [0,577 a 0,782] & 0,625 [0,524 a 0,731] & 0,691 [0,596 a 0,791] \\
    AUC por mama             & 0,832 [0,768 a 0,889] & 0,837 [0,777 a 0,891] & 0,837 [0,767 a 0,901] \\
    Sens. por mama (10\% alertas) & 0,560 [0,407 a 0,700] & 0,600 [0,440 a 0,727] & 0,600 [0,452 a 0,741] \\
    Sens. por mama (20\% alertas) & 0,720 [0,571 a 0,842] & 0,660 [0,531 a 0,796] & 0,700 [0,558 a 0,841] \\
    \bottomrule
  \end{tabular}
  \fonte{Elaborada pelo autor.}
\end{table}

Em todas as métricas, os intervalos de confiança dos três braços se sobrepõem amplamente; não foi possível detectar diferença entre eles com o tamanho da validação. O braço com CLAHE teve as menores estimativas em nível de lesão em todos os pontos da FROC (\autoref{tab:froc_bracos}), mas a diferença está dentro da incerteza. O braço sem recorte teve estimativas próximas às do B0.

\begin{table}[htb]
  \centering
  \caption{Sensibilidade em nível de lesão por FPPI nos três braços (validação).}
  \label{tab:froc_bracos}
  \begin{tabular}{@{}lcccccc@{}}
    \toprule
    Braço & 0,125 & 0,25 & 0,5 & 1 & 2 & 4 \\
    \midrule
    B0 & 0,408 & 0,466 & 0,563 & 0,612 & 0,689 & 0,786 \\
    B2 & 0,359 & 0,427 & 0,485 & 0,563 & 0,641 & 0,728 \\
    B4 & 0,427 & 0,505 & 0,563 & 0,650 & 0,728 & 0,767 \\
    \bottomrule
  \end{tabular}
  \fonte{Elaborada pelo autor.}
\end{table}

\section{Classificador por mama}

\pendente{kappa ponderado do BI-RADS, macro-F1, acurácia, AUC para BI-RADS 4 ou superior e métricas de densidade na validação; matriz de confusão}

\section{Extração de informação dos laudos}

\pendente{concordância entre os dois anotadores (kappa) e desempenho do extrator E0 (e E1, se houver) por campo}

\section{Auditoria}

\pendente{limiar escolhido; desempenho por tipo de erro injetado; comparação com as referências Z0, Z1 e Z2; exemplos de alertas com evidência}

\section{Avaliação final no teste}\label{sec:res_teste}

\pendente{resultados no teste oficial do VinDr-Mammo e no teste externo CBIS-DDSM, com os modelos e limiares congelados, executados uma única vez}

% =============================================================================
\chapter{Discussão}\label{cap:discussao}
% =============================================================================

\section{Desempenho do detector}

O detector congelado encontra cerca de 61\% das massas com um falso positivo por imagem e, em nível de mama, cerca de 56\% das mamas com massa quando 10\% das mamas normais geram alerta. Para o papel que o detector tem neste trabalho, a leitura em nível de mama é a mais relevante: o auditor não precisa delimitar a lesão com precisão, e sim decidir se uma mama tem um achado suspeito que o laudo deveria mencionar.

A comparação com a literatura exige cuidado. O Mammo-CLIP relata mAP de cerca de 0,58 para massas no VinDr-Mammo \cite{ghosh2024mammoclip}, acima do AP de 0,448 obtido aqui nas imagens com achado. Porém, os protocolos não são iguais: este trabalho usa uma validação própria, retirada do treino oficial, e reporta o AP com e sem as imagens normais, que mudam o resultado em mais de duas vezes (\autoref{tab:runs_detector}). Sem o mesmo conjunto de avaliação e a mesma definição de quais imagens entram na conta, a diferença numérica não permite concluir que um modelo é melhor que o outro. \pendente{conferir no artigo do Mammo-CLIP quais imagens entram no cálculo do mAP e ajustar este parágrafo}

O efeito das imagens normais também orienta a escolha das métricas. O AP na validação completa é baixo (0,186) porque a validação tem cerca de doze imagens normais para cada imagem com massa, a mesma proporção esperada na prática. A FROC e as métricas por mama expressam esse custo de forma direta, em falsos positivos por imagem e em fração de mamas normais com alerta, que é o que determina quantos alertas um revisor precisaria descartar.

\section{Efeito do pré-processamento}

A ablação não detectou diferença entre os braços. Dois pontos merecem comentário. O primeiro é o CLAHE: embora seja uma etapa comum, ele teve as menores estimativas em nível de lesão em todos os pontos de operação. A diferença está dentro da incerteza, mas o resultado não apoia o uso do CLAHE como passo padrão para detecção de massas com este detector. Uma hipótese é que o realce local amplifica a textura do tecido fibroglandular, que se parece com o contorno de uma massa, e aumenta os falsos positivos.

O segundo ponto é o recorte da mama. Em classificação no VinDr-Mammo, o recorte trouxe ganho de cerca de quatro pontos de AUC \cite{ibragimov2024mamt4}; aqui, na detecção de massas, o braço sem recorte ficou próximo da referência. Uma explicação possível é que, na resolução usada, a mama já ocupa a maior parte da imagem depois do espelhamento e do redimensionamento, e o ganho de resolução efetiva com o recorte é pequeno. Outra é que um detector, que avalia regiões locais, depende menos da área ocupada pelo fundo do que um classificador global. Os dados deste trabalho não permitem separar essas explicações.

O resultado nulo da ablação também tem valor prático. Como nenhum braço foi melhor, a referência B0 foi mantida, e o recorte continua sendo usado por reduzir o tamanho dos arquivos e o custo de processamento, e não por um ganho de desempenho demonstrado.

\section{Limitações}

\begin{itemize}
  \item \textbf{Tamanho da validação.} A validação tem 50 mamas com massa, o que produz intervalos de confiança largos (cerca de 30 pontos percentuais na sensibilidade por mama). Diferenças de poucos pontos entre variantes não podem ser detectadas.
  \item \textbf{Classes raras.} As classes do detector além de massa têm poucos exemplos e não puderam ser avaliadas. A auditoria de omissão fica, na prática, restrita a massas.
  \item \textbf{Calcificações.} A detecção de calcificações foi retirada do escopo; omissões de calcificações não são cobertas pelo sistema.
  \item \textbf{Gabarito da imagem.} O VinDr-Mammo marca com caixa apenas achados de categoria BI-RADS 3 ou superior \cite{nguyen2023vindr}. Achados benignos presentes na imagem não têm caixa e, se detectados, contam como falsos positivos.
  \item \textbf{Ausência de radiologista.} Não houve revisão clínica das saídas do sistema. As referências são as anotações públicas das bases e as anotações de laudo feitas por anotadores treinados no protocolo de anotação, e a análise de erro limita-se a causas técnicas observáveis.
  \item \textbf{Semente única.} Cada configuração foi treinada uma vez, com semente fixa. A variação entre sementes não foi medida e se soma à incerteza dos intervalos reportados.
\end{itemize}

\pendente{discutir o classificador, o extrator e a auditoria quando os resultados estiverem prontos}

% =============================================================================
\chapter{Conclusão}\label{cap:conclusao}
% =============================================================================

Este trabalho propõe um sistema de auditoria automática de laudos de mamografia, que cruza a imagem com o laudo para sinalizar omissões de achados suspeitos e discordâncias de descritores BI-RADS.

Até o momento, o módulo de imagem está construído e avaliado na validação. O pipeline de pré-processamento parte do DICOM, recorta a mama de forma automática e foi verificado por testes e por inspeção visual, o que levou à correção de uma falha de recorte em mamas pequenas. O detector de massas congelado alcançou AUC de 0,832 em nível de mama e encontra cerca de 56\% das mamas com massa quando 10\% das mamas normais geram alerta. A ablação de pré-processamento não detectou diferença entre usar ou não o realce por CLAHE e entre recortar ou não a mama.

\pendente{completar com os resultados do classificador, do extrator, da auditoria e do teste final; responder a cada objetivo específico do \autoref{cap:introducao}}

\section{Trabalhos futuros}

\begin{itemize}
  \item Incluir a detecção de calcificações, que exige maior resolução e um detector próprio.
  \item Ampliar o conjunto de validação ou usar validação cruzada, para reduzir a incerteza das comparações.
  \item Avaliar o sistema com laudos de outras instituições brasileiras e, com a participação de radiologistas, medir a utilidade clínica dos alertas.
  \item \pendente{completar ao final do trabalho}
\end{itemize}


% ============================================================ PÓS-TEXTUAIS
\postextual
\bibliography{referencias}

\begin{apendicesenv}
\partapendices
\chapter{Checklist CLAIM 2024}
\pendente{preencher os itens do checklist CLAIM \cite{tejani2024claim} à medida que a metodologia for escrita (S16)}
\end{apendicesenv}

\printindex
\end{document}
